%%%%%%%%%%%%%%%%%%%%%%%%%%%%%%%%%%%%%%%%%
%%% WP01
%%%%%%%%%%%%%%%%%%%%%%%%%%%%%%%%%%%%%%%%%

\begin{prtwo}
\tsubsubsection{WP6 - Ethic Requirements}
\end{prtwo}

%%%%%%%%%%%%%%%%%%%%%%%%%%%%%%%%%%%%%%%%%
%%% Section content, please change!
%%%%%%%%%%%%%%%%%%%%%%%%%%%%%%%%%%%%%%%%%

\subsubsection*{Overview and Goals}

\begin{table}[H]
    \renewcommand{\arraystretch}{1.50}		
    \footnotesize   
    {\color{black}\begin{tabular}{*{3}{|p{0.10\textwidth}}|l|}
        \hline
        \rowcolor{mygray} \multicolumn{4}{|c|}{\textit{\color{white}Work Package Summary}} \\
        \hline
        \rowcolor{mylightergray} \textit{WP No.} & \cellcolor{white} 6 & \textit{Title of WP} & \cellcolor{white} Ethics requirements \\
        \hline
        \rowcolor{mylightergray} \textit{Start} & \cellcolor{white} M1 & \textit{End} & \cellcolor{white} M48 \\
        \hline
        \rowcolor{mylightergray} \multicolumn{4}{|p{0.978\textwidth}|}{\textit{Participating Organisations}} \\
        \hline
        \multicolumn{4}{|p{0.978\textwidth}|}{
            \hspace*{-0.75cm} 
            \begin{minipage}[t]{\textwidth}
    			\begin{itemize}
    			    \item WP Leader : INFN
                        \item Ethics Advisor : Prof. (emeritus) M. Harakeh, Rijksuniversiteit Groningen, Netherlands 
    			\end{itemize} 
    			\vspace*{0.10em}
			\end{minipage}
        } \\
        \hline
    \end{tabular}}
    \vspace{0.5em}\vfill
    {\color{black}\begin{tabular}{|p{0.978\textwidth}|}
        \hline
        \rowcolor{mylightergray} \textit{Goals} \\
        \hline
        \rowcolor{white} 
        \hspace*{-0.75cm} 
        \begin{minipage}[t]{\textwidth}
    		\begin{itemize}
    		    \item Setup and ensure compliance with the ethics requirements.
    		\end{itemize} 
    		\vspace*{0.10em}
		\end{minipage}        
        \\
        \hline
    \end{tabular}}
    \vspace{0.5em}\vfill
    \resizebox{\textwidth}{!}{%
    {\color{black}\begin{tabular}{|l|>{\centering\arraybackslash}p{0.77\textwidth}|}
        \hline    
        \rowcolor{mylightergray} \textit{Participant number} & \textit{1} \\
        \hline
        \rowcolor{white} \cellcolor{mylightergray}\textit{Participant short name} & INFN \\
        \hline
        \rowcolor{white} \cellcolor{mylightergray}\textit{PM per participant~\footnotemark} & -  \\
        \hline        
    \end{tabular}}
    }
\end{table}
\footnotetext{the PM figures correspond to the full duration of the project.}


\subsubsection*{Status}

%\begin{prtwo}
EURO-LABS aligned with the recommendations described in the report by the  Ethics Advisor.
% confirming that all consortium activities align with the relevant ethics requirements and that no new ethical issues have emerged. Additional details are provided in \textbf{EURO-LABS Deliverable 6.1}. 
 There were no specific new requests during RP3. 
%\end{prtwo}
%\begin{prtwo}
The EURO-LABS consortium remained vigilant and committed to maintaining the highest ethical standards across all activities. %All activities aligned with the ethics requirements outlined in the project documentation.
 A summary of the measures that have continued to be taken is presented below. 
%\end{prtwo}

%\subsubsection*{Progress per Task}

\subparagraph{Health and Safety During Experiments and Training Activities} \mbox{}


%\begin{prtwo}
EURO-LABS provides Transnational Access (TA) to a wide network of laboratories engaged in cutting-edge research in Accelerator-Based Sciences. These experiments may involve exposure to ionising radiation, posing potential risks to researchers, the public, and the environment. It is worth noting that all EU-based TA facilities comply with EU radiation protection legislation and their respective national laws. Non-EU facilities also adhere to stringent safety regulations. EURO-LABS has implemented and adopted comprehensive procedures to guarantee the health and safety of personnel during all activities.
%\end{prtwo}

%\begin{prtwo}
Mandatory safety training is provided periodically to all employees, and additional procedures are in place for accelerator operations and large installations. Radiation workers receive regular training and testing. Visitors are required to complete safety training and maintain health insurance. All users are equipped with personal dosimeters.
%\end{prtwo}

%\begin{prtwo}
During schools and hands-on training sessions organized by EURO-LABS, the same safety protocols—particularly those concerning ionising radiation, in case of activities implying exposure—are applied to visiting lecturers and students.
%\end{prtwo}

\subparagraph{Development of Machine Learning Methods} \mbox{}

%\begin{prtwo}
Work Package 5.3 involves the development and application of Machine Learning (ML) methods to enhance beam quality, transport efficiency, and reproducibility, thereby reducing the time required for beam tuning.
%\end{prtwo}

%\begin{prtwo}
To ensure ethical compliance, emphasis is placed on transparency, robustness, and avoiding bias in ML methods. Although the techniques could have applications in the medical sector, this is outside the scope of the current project.
%\end{prtwo}

%\begin{prtwo}
The ML development adheres to the \textit{Ethics by Design} and \textit{Ethics of Use} principles, as outlined in the European Commission documents dated 25 November 2021. These principles are followed throughout all phases: design, acquisition, implementation, and monitoring.
%\end{prtwo}

%\begin{prtwo}
The goal of this ML project is to reduce the time required for tuning the beams provided to Transnational Users of EURO-LABS. This will help increase and optimize the time available for research and technological studies, while also reducing power consumption.
The ML applications are restricted to a limited set of machine parameters affecting beam optics and cannot modify arbitrary system settings. Any algorithmic bias could only impact quantifiable properties of the beam (e.g., intensity or size), which are continuously monitored.
%\end{prtwo}

%\begin{prtwo}
The ML algorithms used are fully deterministic and robust against minor external variations. Their application in the project is focused, reliable, transparent, and free of ethical concerns, aimed solely at improving the efficiency of beam tuning for Transnational Users.
%\end{prtwo}

% \subsubsection*{Main Results and Achievements}

\subsubsection*{Deviations and Corrective Actions}

%\begin{prtwo}
None
%\end{prtwo}

\subsubsection*{Milestones and Deliverables}

None

%\begin{prtwo}
%{\fontsize{9}{11}\selectfont
%\begin{center}
  %{\color{black}\begin{tabular}[t]{!{\color{mygray}\vrule}p{0.10\linewidth}!
  %{\color{mygray}\vrule}p{0.60\linewidth}!
  %{\color{mygray}\vrule}p{0.20\linewidth}!{\color{mygray}\vrule} } \hline
    %\rowcolor{mycyan} & {\bf Title} & {\bf Status} \\ \hline
%    \cellcolor{mycyan}{\bf D6.1}: & Ethics - OEI - 24m Update & Achieved \\ \hline
%  \end{tabular}}
%\end{center}
%}
%\end{prtwo}

\subsubsection*{Project Meetings}

%\begin{prtwo}
No meetings
%\end{prtwo}
%  {\color{mygray}\vrule}p{0.40\linewidth}!

%%%%%%%%%%%%%%%%%%%%%%%%%%%%%%%%%%%%%%%%%
\begin{prtwo}
\tsubsubsection{Additional Requirements}
\end{prtwo}

The following points, reported in italics, refer to requirements after the completion of the project. They are followed by our responses. 

1. SUSTAINABILITY OF WEBSITE: {\it Make sure that the project website is sustainable for the next years, as a minimum in a static form, but of course preferably with new added content. }

The EURO-LABS website will be maintained by INFN for several years and will be accessible and available in its current form. It will be updated with new results if and when they will be available. 

2. PROTECTION OF RESULTS: {\it If the conditions of the GA are met, adequately protect the results for an appropriate period and with appropriate territorial coverage (Article 27). }

We confirm to have been informed.

3. DISSEMINATION \& EXPLOITATION OF RESULTS: {\it Disseminate without delay and exploit the project results, unless it goes against legitimate interests. Exploitation should be pursued up to four years after the end of the action (Articles 28 and 29.1 Grant Agreement). Explore opportunities, services, tools and guidance developed by the EC: Horizon Results Booster to apply during and after the end of the grant to get assistance in your D\&E activities. Horizon Results Platform to publish key exploitable results for visibility. Horizon Impact Award that targets both FP7 and H2020 completed projects. Consider participation in cluster events. }

We confirm to have been informed.

4. VISIBILITY OF EU FUNDING: {\it When exploiting, protecting, disseminating results (in any form, even when combined with other data) or when undertaking communication activities promoting the action, include the reference to EU funding set out in the GA. (Articles 27, 28, 29 and 38 Grant Agreement) }

We have always made reference to the EU funding, as set out in the GA, in presentations, publications and newsletters. We will continue to do this also after the end of the project for an appropriate period.

5. DISCLAIMER EXCLUDING THE AGENCY’S RESPONSIBILITY: {\it Do not forget the disclaimer excluding the Agency responsibility for any communication activity and any dissemination of results (Articles 38.1.3 and Art 29.5)
 }

We confirm to have been informed.

6. OPEN ACCESS: {\it Deposit all peer-reviewed publications in a repository and give Open Access as soon as possible, at the latest (6 or 12) months after publication (Article 29.2). }

We confirm to have been informed.

7. REPORTING ON PUBLICATIONS: {\it Continue to report all publications via your project's "continuous reporting module". This will allow us to have accurate and realistic data on our funded projects. }

We confirm to have been informed.

8. KEEPING RECORDS: {\it Keep records and other supporting documentation (original documents) for a period of five years after the final payment (Article 18). }

We confirm to have been informed.

9. PROVIDING INFORMATION: {\it Upon request, provide any information requested in order to verify eligibility of the costs, proper implementation of the action and compliance with any other obligation under the grant agreement (Article 17). In addition, any information relevant to evaluate the impact of the action can be requested up to five years after the final payment (Article 23.1).}

We confirm to have been informed.

